\documentclass{article}
\usepackage{amsmath, amssymb, amsthm}
\usepackage{hyperref}
\usepackage{geometry}
\geometry{margin=1in}

\title{Erd\H{o}s Problem \#953}
\date{Last updated: 2025-08-31}
\author{Source: erdosproblems.com}

\begin{document}
\maketitle

\noindent\textbf{Status:} OPEN \quad \textbf{No prize}

\noindent\textbf{Tags:} geometry, distances

\medskip

\noindent\textbf{Problem Statement:}

Let $A\subset \{ x\in \mathbb{R}^2 : \lvert x\rvert &#60;r\}$ be a measurable set with no integer distances, that is, such that $\lvert a-b\rvert \not\in \mathbb{Z}$ for any distinct $a,b\in A$. How large can the measure of $A$ be?




A problem of Erd\H{o}s and S\'{a}rk\"{o}zi. Erd\H{o}s \cite{Er77c} writes that 'S\'{a}rk\"{o}zi has the sharpest results, but nothing has been published yet'. The trivial upper bound is $O(r)$. Koizumi and Kovac have observed in the comments that S\'{a}rk\"{o}zy's lower bound in [466] can be adapted to give a lower bound of $\gg_\epsilon r^{1/2-\epsilon}$ for all $\epsilon&gt;0$. See also [465] for a similar problem (concerning upper bounds) and [466] for a similar problem (concerning lower bounds).




References

[Er77c] Erd\H{o}s, Paul, Problems and results on combinatorial number theory. III . Number theory day (Proc. Conf., Rockefeller Univ., New York, 1976) (1977), 43-72.


\bigskip
\noindent\small{Source: \url{https://www.erdosproblems.com/953}}
\end{document}
